\scanpage{381}
\mspage{120}
% archive sheet number: 380
(Sans avoir pour autant à calculer le groupe fondamental de $\underline{\mathrm{Hom}}\ \underline{\mathrm{cof}}(\mathbb N^\circ,T)$ !) Si on a trois éléments $\underline t,\underline t',\underline t''$ de $\underline{\mathrm{Hom}}\ \underline{\mathrm{cof}}$, il suffit de regarder le diagramme déduit du diagramme commutatif suivant
\diagram{\begin{tikzcd}[ampersand replacement=\&,row sep=small,column sep=small]
\&\&t*t'*t''\arrow[dl]\arrow[dr]\&\&\\
\&\underline t*\underline t'\arrow[dl]\arrow[dr]\&\&t'*t''\arrow[dl]\arrow[dr]\&\\
\underline t\&\&\underline t'\&\&\underline t''
\end{tikzcd}}
en lui appliquant le contrafoncteur $\underline s\mathrel{\unsure{\mapstochar\relbar|\mapsto}}C_\infty(\underline s)$ de $\underline{\mathrm{Hom\,cof}}$ dans $\mathrm{Hot}$.

Donc dans le cas \deleted{des} \inserted{$C(a,b)={}$}$\underline{\mathrm{Ch}}_{\unsure{\Sigma|Z}}(X;a,b)$ \deleted{\unsure{$\underline{\mathrm{Ch}}_\Sigma(X)$}}, je n'ai plus aucun scrupule à considérer la catégorie $C_\infty(a,b)$ comme un objet \uline{canonique} de $\mathrm{Hot}$, relié à \deleted{\unsure{un|au}} l'ind-objet canonique $H_T(a,b)$ \inserted{de $\mathrm{Hot}$} par un lien canonique
\[
C_\infty(a,b)\mathrel{\leftarrow\deleted{\unsure{>}}}H_T(a,b),
\tag*{$(\unsure{\alpha|a})$}
\]
qui dans $\mathrm{Hot}_0^{\unsure{\wedge|1}}$ donne un isomorphisme
\[
C_\infty(a,b)\simeq\underline{\lim}\,H_T(a,b),
\tag*{$(\unsure{\alpha'|\alpha_0|\alpha})$}
\]
i.e. la flèche $\alpha$, induit, pour $\unsure{\zeta|\xi}\in\underline{\mathrm{Ob}}\,\mathrm{Hot}$, une \deleted{\unsure{bije}}application ensembliste
\[
\mathrm{Hom}_{\mathrm{Hot}}\bigl(\zeta,C_\infty(a,b)\bigr)
\mathrel{\leftarrow\deleted{\unsure{>}}}
\underset{\substack{\Vert\\
\underrightarrow{\lim}\,\mathrm{Hom}_{\mathrm{Hot}}(\zeta,C_n(a,b))}}
{\mathrm{Hom}_{\mathrm{Ind\,Hot}}\bigl(\zeta,H_T(a,b)\bigr)}
\]
qui est \uline{bijective} si $\zeta$ est dans $\unsure{\widetilde{\mathrm{Hot}}_0|\mathrm{Hot}_0}$.
